\section{Related Work}
\label{sec:related}

\subsection{Process-Oriented CBR Foundations}

The foundational CBR cycle was given its canonical formulation by Aamodt \& Plaza
\cite{aamodt1994}. Bergmann \cite{bergmann2002} provides foundational work on applying
CBR principles to workflow and experience management; Minor, Montani \& Recio-Garcia
\cite{minor2014} characterize POCBR as the field in which cases capture workflow
knowledge.

Within POCBR, Bergmann \& Gil \cite{bergmann2014} address efficient retrieval through
structural similarity; Muller \& Bergmann \cite{muller2015a} introduce a query language
where the query itself is a partial workflow. Bottrighi et al.\ \cite{bottrighi2016}
extend this to business process operational support. Leake \& Kendall-Morwick
\cite{leake2008} explicitly propose mining workflow provenance traces as CBR cases.
NormCode's contribution to this line of work is to provide a structural guarantee: by
enforcing scope isolation at the language level, every captured checkpoint is a
reliable, self-contained case rather than a trace that may embed implicit shared state.

\subsection{Multi-Level and Hierarchical CBR}

Muller \& Bergmann \cite{muller2015b} introduce generalized workflow cases that cover a
subspace of the case representation space --- the closest antecedent we are aware of to
NormCode's Level~2: both treat the abstract workflow pattern as a case object. The
generative difference is that NormCode's Level~2 cases are fully executable compiled
artifacts and the distillation process is itself expressible as a NormCode plan. Veloso
\cite{veloso1997} and Ihrig \& Kambhampati \cite{ihrig1997} establish that plan
derivations can be stored and retrieved as cases.

\subsection{CBR and Large Language Models}

Rubin, Herzig \& Berant \cite{rubin2022} learn to retrieve few-shot prompts; Wilkerson
\& Leake \cite{wilkerson2024} implement CBR components using LLMs; Lenz, Hoffmann \&
Bergmann \cite{lenz2025} evaluate LLMs as similarity assessors (LLsiM). These works use
CBR to improve individual LLM calls, or LLMs to implement CBR functions.

NormCode addresses a distinct problem: applying a CBR architecture to manage the
\textit{execution of multi-step LLM workflows}. Among prior work, Guo et al.\
\cite{guo2024} (DS-Agent) is the most directly related: CBR is used to iteratively
revise multi-step data-science plans. The structural difference is that DS-Agent's case
objects are narrative plans without enforcement guarantees, whereas NormCode's
checkpoints are structurally isolated by compiler and runtime --- which is the
precondition for reliable CBR retrieval and revision (Sect.~\ref{sec:background}).

\subsection{LLM Workflow Tools}

LangGraph provides checkpoint support and time-travel replay. PromptFlow and LangFlow
provide graph-based visual workflow authoring. AutoGen \cite{wu2023} supports
conversational multi-agent orchestration. These frameworks execute steps by passing
information through accumulated conversation histories or prompt templates that
implicitly reference prior outputs. Relative to NormCode's scope rule, none provides
explicit structural isolation guarantees at the language level: checkpoints may embed
implicit shared state from earlier steps. As a consequence, retrieving and replaying
such a checkpoint in a new execution context may produce incorrect behavior, failures
are non-localizable, and case reuse is not guaranteed to reproduce the original
data-dependency behavior --- the three conditions identified in
Sect.~\ref{sec:background} as preconditions for reliable CBR operations.

\subsection{LLM Agents and the ReAct Pattern}

Yao et al.\ \cite{yao2023} propose ReAct --- interleaving reasoning traces with
grounded action execution in a tight loop. The observe-reason-act pattern has since
become a standard building block for tool-using LLM systems. Reflexion
\cite{shinn2023} extends it with verbal reinforcement: agents verbalize failure
feedback into episodic memory and retry, improving without weight updates. MetaGPT
\cite{hong2024} encodes software-engineering standard operating procedures as prompt
sequences for multi-role agent collaboration; SWE-agent \cite{yang2024} introduces an
agent-computer interface for autonomous software engineering. AutoGen \cite{wu2023}
generalizes the pattern to conversational multi-agent coordination.

NormCode plans that use Canvas Integration facilities implement the ReAct pattern as
explicitly compiled plans. The Canvas Assistant's per-iteration loop --- read user
input $\rightarrow$ classify command $\rightarrow$ execute via \texttt{me.hands} or
\texttt{me.mind} $\rightarrow$ generate response --- is a ReAct cycle expressed as a
structured plan with declared scopes. The Code Assistant's pipeline similarly reasons
over retrieved context and acts via tool calls across seven stages.

Deployed agentic coding tools --- Cursor \cite{cursor2024}, Claude Code
\cite{claudecode2026}, and OpenClaw \cite{openclaw2026} --- demonstrate this pattern
at production scale: each implements a persistent observe-reason-act loop over
codebases, terminals, and file systems. The core workflow of such tools --- triage the
task, explore context, generate and apply changes, verify --- is structurally similar
to the Code Assistant NormCode plan. Expressed as a NormCode Level~2 case, this
workflow gains auditable structure, pre-execution review (C2), and a Level~1 case base
at every agent step. The NormCode framing is therefore not specific to the four plans
presented here; it applies to the class of deployed agentic tools at large.

The structural difference is that in standard ReAct and its derivatives, the
reasoning-action chain is an implicit continuation; the reasoning trace and runtime
state are not formally separated. In NormCode, each step in the agent loop is a node
with a declared scope, a flow index, and an automatic checkpoint. The agent's execution
trace is a Level~1 case base: every iteration step is inspectable, revisable, and
reusable under the CBR operations C1 and C3. The absence of this property in existing
frameworks is independently confirmed by recent work: AgentStepper \cite{agentstepper2026}
introduces post-hoc interactive debugging (breakpoints, step-through, online prompt
editing) for software development agents precisely because the underlying frameworks
provide no structural transparency. NormCode makes these properties structural rather
than retrofitted.

\subsection{Process Mining}

Process mining \cite{vandermaaist2011} discovers process models from event logs.
NormCode's Level~2 distillation is related in spirit to guided process model discovery.
The key difference is authorial intent: process mining performs automated discovery from
logs, while NormCode's compilation is human-directed and LLM-assisted, producing
executable artifacts rather than descriptive models.
